% !TeX program = lualatex
% halo:begin
% {
%   "version": 1,
%   "publish": true,
%   "course": "mathematical-analysis-i",
%   "section": "4.2",
%   "title": "数学分析（一）4.2：连续函数的性质",
%   "categories": ["study-notes"],
%   "tags": ["mathematics"],
%   "excerpt": "连续函数的局部性质、闭区间上的基本定理、反函数的连续性、一致连续性和初等函数的连续性。"
% }
% halo:end
% 来源：IMG_3107—IMG_3111.HEIC；华东师大《数学分析》第5版上册第四章§2—§3。
\RequirePackage{iftex}
\RequireLuaTeX
\documentclass[UTF8,fontset=fandol,a4paper,12pt]{ctexart}
\usepackage[margin=2.5cm]{geometry}
\usepackage{amsmath,amssymb,amsthm}
\usepackage{enumitem,fancyhdr}
\usepackage[hidelinks]{hyperref}
\newcommand{\R}{\mathbb R}
\theoremstyle{definition}
\newtheorem{definition}{定义}[section]
\newtheorem{theorem}[definition]{定理}
\newtheorem{proposition}[definition]{命题}
\newtheorem{example}[definition]{例题}
\renewcommand{\proofname}{证明}
\setlist[enumerate]{leftmargin=2em,itemsep=0.3em,topsep=0.3em}
\setlength{\headheight}{16pt}
\pagestyle{fancy}\fancyhf{}
\fancyhead[L]{\small 数学分析（一）}
\fancyhead[R]{\small 4.2\quad 连续函数的性质}
\fancyfoot[C]{\thepage}\renewcommand{\headrulewidth}{0.3pt}
\raggedbottom
\title{连续函数的性质}
\author{Yuchen Zhang}
\date{整理日期：2026年11月6日}
\makeatletter
\renewcommand{\maketitle}{\begin{center}{\LARGE\bfseries\@title\par}\vspace{0.6em}\@author\par\vspace{0.2em}{\small\@date\par}\end{center}}
\makeatother
\begin{document}
\maketitle
\section{连续函数的局部性质}
函数在一点连续，便能用这一点的函数值控制附近的函数值。先看由此得到的局部有界性和局部保号性。
\begin{theorem}[局部有界性与局部保号性]
设 $f$ 在 $x_0$ 连续，则：
\begin{enumerate}
\item 存在 $\delta>0$ 和 $M>0$，使 $|x-x_0|<\delta$ 时有 $|f(x)|<M$。
\item 若 $f(x_0)>0$，则在 $x_0$ 的某个邻域内 $f(x)>0$；若 $f(x_0)<0$，则在某个邻域内 $f(x)<0$。
\end{enumerate}
\end{theorem}
\begin{proof}
取 $\varepsilon=1$，由连续性可取 $\delta>0$，使 $|x-x_0|<\delta$ 时
\[
 |f(x)|\le |f(x_0)|+|f(x)-f(x_0)|<|f(x_0)|+1.
\]
这就证明了局部有界性。若 $f(x_0)>0$，取 $\varepsilon=f(x_0)/2$，则在充分小的邻域内
\[
 f(x)>f(x_0)-\frac{f(x_0)}2=\frac{f(x_0)}2>0.
\]
负值情形对 $-f$ 应用同样结论即可。在区间端点，上述邻域均限制在定义区间内。
\end{proof}
\begin{theorem}[连续函数的运算]
若 $f,g$ 在 $x_0$ 连续，则 $f\pm g$、$fg$ 在 $x_0$ 连续；若 $g(x_0)\ne0$，则 $f/g$ 也在 $x_0$ 连续。若 $g$ 在 $x_0$ 连续、$f$ 在 $u_0=g(x_0)$ 连续，且 $f(g(x))$ 在 $x_0$ 附近有定义，则复合函数 $f\circ g$ 在 $x_0$ 连续。
\end{theorem}
\begin{proof}
四则运算的结论由极限的运算法则得到；分母在 $x_0$ 附近非零由局部保号性保证。对复合函数，给定 $\varepsilon>0$，由 $f$ 在 $u_0$ 连续可取 $\eta>0$，使 $|u-u_0|<\eta$ 时 $|f(u)-f(u_0)|<\varepsilon$。再由 $g$ 连续取 $\delta>0$，使 $|x-x_0|<\delta$ 时 $|g(x)-u_0|<\eta$，便得所需结论。
\end{proof}
\section{闭区间上的有界性与最值}
局部有界性只控制一个点附近的函数值。若函数在整个闭区间连续，还能控制全区间，并保证最大值和最小值确实取到。
以下用 $f\in C[a,b]$ 表示 $f$ 在 $[a,b]$ 上连续，其中端点处按单侧连续理解。
\begin{theorem}[有界性定理]
若 $f\in C[a,b]$，则 $f$ 在 $[a,b]$ 上有界。
\end{theorem}
\begin{proof}
反设 $f$ 无界，则对每个正整数 $n$，可取 $x_n\in[a,b]$，使 $|f(x_n)|>n$。点列 $\{x_n\}$ 有界，由致密性定理存在子列 $x_{n_k}\to x_0$。因 $a\le x_{n_k}\le b$，有 $x_0\in[a,b]$。由连续性，$f(x_{n_k})\to f(x_0)$，故该函数值列有界；但 $|f(x_{n_k})|>n_k\to+\infty$，矛盾。
\end{proof}
\begin{theorem}[最大值与最小值定理]
若 $f\in C[a,b]$，则存在 $x_m,x_M\in[a,b]$，使对一切 $x\in[a,b]$ 都有
\[
 f(x_m)\le f(x)\le f(x_M).
\]
\end{theorem}
\begin{proof}
由有界性定理，$M=\sup\{f(x):x\in[a,b]\}$ 为有限实数。对每个正整数 $n$，由上确界的定义可取 $x_n\in[a,b]$，使
\[
 M-\frac1n<f(x_n)\le M.
\]
故 $f(x_n)\to M$。再从有界点列 $\{x_n\}$ 中取子列 $x_{n_k}\to x_M\in[a,b]$。由连续性及极限的唯一性，
\[
 f(x_M)=\lim_{k\to\infty}f(x_{n_k})=M.
\]
因此上确界能够取到，即为最大值。对 $-f$ 应用这一结论，便得到 $f$ 的最小值。
\end{proof}
闭区间条件在这里保证极限点仍属于定义域。例如 $f(x)=x$ 在 $(0,1)$ 上连续且有界，却既无最大值也无最小值；$f(x)=1/x$ 在 $(0,1)$ 上连续，但无界。
\section{零点定理与介值定理}
连续函数在两端异号时，必在中间取到零值。
\begin{theorem}[零点定理]
若 $f\in C[a,b]$ 且 $f(a)f(b)<0$，则存在 $c\in(a,b)$，使 $f(c)=0$。
\end{theorem}
\begin{proof}
不妨设 $f(a)<0<f(b)$。令
\[
 S=\{x\in[a,b]:f(x)<0\},\qquad c=\sup S.
\]
由于 $a\in S$ 且 $S\subseteq[a,b]$，上确界存在。由两端的连续性与保号性，$a$ 右侧有属于 $S$ 的点，而 $b$ 左侧的一小段不含 $S$ 中的点，故 $a<c<b$。

由上确界的定义可取 $x_n\in S$，使 $c-1/n<x_n\le c$，于是 $x_n\to c$。由连续性，
\[
 f(c)=\lim_{n\to\infty}f(x_n)\le0.
\]
若 $f(c)<0$，则由保号性，$c$ 右侧充分近的点也有负函数值，从而属于 $S$，与 $c$ 是 $S$ 的上界矛盾。因此 $f(c)=0$。
\end{proof}
零点定理也能用来研究函数是否取到一个给定值，只需从函数值中减去这个数。
\begin{theorem}[介值定理]
若 $f\in C[a,b]$，$\mu$ 严格介于 $f(a)$ 与 $f(b)$ 之间，则存在 $c\in(a,b)$，使 $f(c)=\mu$。
\end{theorem}
\begin{proof}
令 $F(x)=f(x)-\mu$。则 $F\in C[a,b]$，且 $F(a)F(b)<0$。由零点定理存在 $c\in(a,b)$，使 $F(c)=0$，即 $f(c)=\mu$。
\end{proof}
若 $m,M$ 分别是 $f$ 在 $[a,b]$ 上的最小值与最大值，则在取到最值的两个点之间应用介值定理，得到
\[
 \boxed{f([a,b])=[m,M].}
\]
更一般地，连续函数把一个区间映成一个区间；因为任意两个函数值之间的数，都可在对应自变量之间取到。
\begin{example}
设 $f\in C[a,b]$，且 $f(a)<a$、$f(b)>b$。证明存在 $c\in(a,b)$，使 $f(c)=c$。
\end{example}
\begin{proof}
令 $F(x)=f(x)-x$，则 $F\in C[a,b]$，且 $F(a)<0<F(b)$。由零点定理，存在 $c\in(a,b)$，使 $F(c)=0$，即 $f(c)=c$。
\end{proof}
\section{反函数的连续性}
设 $f$ 在区间 $I$ 上连续，则值域 $J=f(I)$ 也是区间。若 $f$ 还严格单调，同一函数值便只对应一个自变量，因而存在反函数 $g:J\to I$。
\begin{theorem}[反函数的连续性]
区间上的连续严格单调函数，其反函数也是连续严格单调函数，且单调方向相同。
\end{theorem}
\begin{proof}
只证 $f$ 严格递增的情形。由严格递增性可直接得到 $g$ 严格递增。任取 $y_0=f(x_0)\in J$，先设 $x_0$ 为 $I$ 的内点。给定 $\varepsilon>0$，在 $I$ 中选取 $x_-,x_+$，使
\[
 x_0-\varepsilon<x_-<x_0<x_+<x_0+\varepsilon.
\]
由严格递增性，$f(x_-)<y_0<f(x_+)$。取
\[
 \delta=\min\{y_0-f(x_-),\ f(x_+)-y_0\}>0.
\]
当 $y\in J$ 且 $|y-y_0|<\delta$ 时，有 $f(x_-)<y<f(x_+)$，所以
\[
 x_-<g(y)<x_+,\qquad |g(y)-g(y_0)|<\varepsilon.
\]
故 $g$ 在 $y_0$ 连续。若 $x_0$ 是属于 $I$ 的端点，只从区间内选取一个界点，便得到相应的单侧连续性。严格递减的情形同理。
\end{proof}
\section{一致连续性}
在一点连续时，$\delta$ 可以随考察的点改变；一致连续要求同一个 $\delta$ 控制定义域内任意一对足够接近的点。
\begin{definition}
设 $f$ 在区间 $I$ 上有定义。若对任意 $\varepsilon>0$，存在 $\delta>0$，使对任意 $x,y\in I$，只要 $|x-y|<\delta$，就有 $|f(x)-f(y)|<\varepsilon$，则称 $f$ 在 $I$ 上一致连续。
\end{definition}
这里的 $\delta$ 只依赖于 $\varepsilon$，不依赖于 $x,y$ 的位置。固定 $y=x_0$ 即知，一致连续必推出逐点连续；反过来一般不成立。

一致连续的否定形式是：存在 $\varepsilon_0>0$，对每个 $\delta>0$，都能找到 $x,y\in I$，满足
\[
 |x-y|<\delta,\qquad |f(x)-f(y)|\ge\varepsilon_0.
\]
\begin{theorem}[一致连续的数列判别]
$f$ 在 $I$ 上一致连续，当且仅当对任意两列 $x_n,y_n\in I$，都有
\[
 |x_n-y_n|\to0\quad\Longrightarrow\quad
 |f(x_n)-f(y_n)|\to0.
\]
\end{theorem}
\begin{proof}
若 $f$ 一致连续，给定 $\varepsilon>0$，取定义中的 $\delta$。充分大的 $n$ 满足 $|x_n-y_n|<\delta$，因而函数值之差小于 $\varepsilon$。

反之，若 $f$ 不一致连续，由否定形式对 $\delta=1/n$ 取点 $x_n,y_n$，使
\[
 |x_n-y_n|<\frac1n,\qquad |f(x_n)-f(y_n)|\ge\varepsilon_0.
\]
这与数列条件矛盾。
\end{proof}
\begin{example}
证明 $\sin x$ 在 $\R$ 上一致连续。
\end{example}
\begin{proof}
由和差化积公式及 $|\sin t|\le|t|$，
\[
 |\sin x-\sin y|
 =2\left|\sin\frac{x-y}{2}\cos\frac{x+y}{2}\right|
 \le |x-y|.
\]
给定 $\varepsilon>0$，取 $\delta=\varepsilon$ 即可。
\end{proof}
\begin{example}
证明 $1/x$ 在 $[d,1]$（$0<d<1$）上一致连续，但在 $(0,1)$ 上不一致连续。
\end{example}
\begin{proof}
若 $x,y\in[d,1]$，则
\[
 \left|\frac1x-\frac1y\right|=\frac{|x-y|}{xy}\le\frac{|x-y|}{d^2}.
\]
取 $\delta=d^2\varepsilon$，便得一致连续性。

在 $(0,1)$ 中取 $x_n=1/(n+1)$、$y_n=1/(n+2)$，则
\[
 |x_n-y_n|=\frac1{(n+1)(n+2)}\to0,
 \qquad \left|\frac1{x_n}-\frac1{y_n}\right|=1.
\]
由数列判别，$1/x$ 在 $(0,1)$ 上不一致连续。
\end{proof}
\begin{example}
证明 $\tan x$ 在 $[0,\pi/2)$ 上不一致连续。
\end{example}
\begin{proof}
令 $t_n=\pi/2^{n+1}$，取 $x_n=\pi/2-t_n$、$y_n=\pi/2-t_n/2$。则 $x_n,y_n\in[0,\pi/2)$，且 $|y_n-x_n|=t_n/2\to0$。由三角恒等式，
\[
 \tan y_n-\tan x_n
 =\cot\frac{t_n}{2}-\cot t_n
 =\frac1{\sin t_n}\ge1.
\]
因此函数值之差不趋于零，$\tan x$ 不一致连续。
\end{proof}
后两个例子中，两列点趋近端点时，自变量之差趋于零，函数值之差却不趋于零。闭区间上的连续性则足以排除这种情形。
\begin{theorem}[闭区间上的一致连续性定理]
若 $f\in C[a,b]$，则 $f$ 在 $[a,b]$ 上一致连续。
\end{theorem}
\begin{proof}
反设不一致连续，则存在 $\varepsilon_0>0$ 及 $x_n,y_n\in[a,b]$，使
\[
 |x_n-y_n|<\frac1n,\qquad |f(x_n)-f(y_n)|\ge\varepsilon_0.
\]
由致密性定理，存在子列 $x_{n_k}\to x_0\in[a,b]$。又由
\[
 |y_{n_k}-x_0|\le|y_{n_k}-x_{n_k}|+|x_{n_k}-x_0|\to0,
\]
得 $y_{n_k}\to x_0$。由 $f$ 在 $x_0$ 连续，$f(x_{n_k})$ 与 $f(y_{n_k})$ 都趋于 $f(x_0)$，故两者之差趋于零，与其绝对值不小于 $\varepsilon_0$ 矛盾。
\end{proof}
除了直接利用闭区间，还可以把相邻两段上的一致连续性拼接起来；关键是控制分处接点两侧的两个点。
\begin{proposition}[相邻区间的拼接]
设区间 $I_1,I_2$ 满足 $\sup I_1=\inf I_2=c$ 且 $c\in I_1\cap I_2$。若同一函数 $f$ 在 $I_1,I_2$ 上分别一致连续，则它在 $I_1\cup I_2$ 上一致连续。
\end{proposition}
\begin{proof}
给定 $\varepsilon>0$，分别对 $\varepsilon/2$ 取一致连续定义中的 $\delta_1,\delta_2$，令 $\delta=\min\{\delta_1,\delta_2\}$。任取 $x,y\in I_1\cup I_2$ 且 $|x-y|<\delta$。若它们属于同一区间，函数值之差小于 $\varepsilon/2$。若分处两侧，不妨设 $x\le c\le y$，则
\[
 |x-c|\le|x-y|<\delta,\qquad |y-c|\le|x-y|<\delta,
\]
从而
\[
 |f(x)-f(y)|\le|f(x)-f(c)|+|f(c)-f(y)|<\frac\varepsilon2+\frac\varepsilon2=\varepsilon.
\]
\end{proof}
\section{初等函数的连续性}
前面的四则运算和复合函数连续性定理，使我们能由基本初等函数的连续性得到一般初等函数的连续性。
\begin{theorem}
基本初等函数在各自的定义区间上连续：指数函数 $a^x$（$a>0$，$a\ne1$）在 $\R$ 上连续，对数函数 $\log_a x$ 在 $(0,+\infty)$ 上连续；幂函数、三角函数与反三角函数在各自的定义区间上连续。
\end{theorem}
\begin{proof}
由前面已证的极限 $\lim_{h\to0}e^h=1$，对任意 $x_0\in\R$，
\[
 \lim_{x\to x_0}e^x=e^{x_0}\lim_{x\to x_0}e^{x-x_0}=e^{x_0}.
\]
因此 $\exp x=e^x$ 在 $\R$ 上连续。对固定的 $a>0$，由 $a^x=\exp(x\ln a)$ 及复合连续性，得到一般指数函数的连续性；这里的 $\ln a$ 是固定常数。对数函数是连续严格单调指数函数的反函数，故连续。

对 $x>0$，由 $x^\alpha=\exp(\alpha\ln x)$ 及复合连续性可得幂函数连续。整数幂的连续性由四则运算得到，开方函数的连续性由反函数连续性得到；在负半轴有定义的有理幂，可通过 $x\mapsto-x$ 归结到正半轴。若 $\alpha>0$ 且 $0$ 属于定义域，则对任意 $\varepsilon>0$，取 $\delta=\varepsilon^{1/\alpha}$，由正幂函数的单调性，$|x|<\delta$ 时有 $|x^\alpha|=|x|^\alpha<\varepsilon$。故在 $0$ 处也连续；按定义域的实际情形，取单侧或双侧连续。

由 $|\sin x-\sin y|\le|x-y|$ 可得正弦函数连续，余弦函数由 $\cos x=\sin(x+\pi/2)$ 得到。正切、余切在分母非零的区间连续。将正弦、余弦、正切、余切分别限制在其通常的严格单调区间，再用反函数连续性定理，得到相应反三角函数的连续性。
\end{proof}
\begin{theorem}
任何初等函数在其定义区间上连续。
\end{theorem}
\begin{proof}
初等函数由基本初等函数经过有限次四则运算和复合而成，逐次应用连续函数的运算定理即可。涉及除法时，所讨论的区间须使分母非零。
\end{proof}
例如 $\sqrt{1-x^2}$ 在 $[-1,1]$ 上连续，其中在 $-1$ 处取右连续、在 $1$ 处取左连续，与前面对区间端点的约定一致。
\end{document}
